\begin{lemma}
\label{lemma-tricanonical}
Let $k$ be a field. Let $X$ be a proper scheme over $k$ of dimension $1$ with
$H^0(X, \mathcal{O}_X) = k$ having genus $g \geq 2$. Assume the singularities
of $X$ are at-worst-nodal and $\omega_X$ is ample. Then
$\omega_X^{\otimes 3}$ is very ample and $H^1(X, \omega_X^{\otimes 3}) = 0$.
\end{lemma}

\begin{proof}
Combining Varieties, Lemma
\ref{varieties-lemma-ampleness-in-terms-of-degrees-components} and
Lemmas \ref{lemma-rational-tail-negative} and \ref{lemma-rational-bridge-zero}
we see that $X$ contains no rational tails or bridges.
Then we see that $\omega_X^{\otimes 3}$ is globally generated
by Lemma \ref{lemma-contracting-rational-tails}.
Choose a $k$-basis $s_0, \ldots, s_n$ of
$H^0(X, \omega_X^{\otimes 3})$. We get a morphism
$$
\varphi_{\omega_X^{\otimes 3}, (s_0, \ldots, s_n)} :
X \longrightarrow \mathbf{P}^n_k
$$
See Constructions, Section \ref{constructions-section-projective-space}.
The lemma asserts that this morphism is a closed immersion.
To check this we may replace $k$ by its algebraic closure, see
Descent, Lemma \ref{descent-lemma-descending-property-closed-immersion}.
Thus we may assume $k$ is algebraically closed.

\medskip\noindent
Assume $k$ is algebraically closed.
We will use Varieties, Lemma
\ref{varieties-lemma-variant-separate-points-tangent-vectors}
to prove the lemma.
Let $Z \subset X$ be a closed subscheme of degree $2$ over $Z$
with ideal sheaf $\mathcal{I} \subset \mathcal{O}_X$.
We have to show that
$$
H^0(X, \mathcal{L}) \to H^0(Z, \mathcal{L}|_Z)
$$
is surjective. Thus it suffices to show that
$H^1(X, \mathcal{I}\mathcal{L}) = 0$.
To do this we will use Lemma \ref{lemma-vanishing-on-gorenstein}.
Thus it suffices to show that
$$
3\deg(\omega_X|_Y) > -2\chi(Y, \mathcal{O}_Y) + \deg(Z \cap Y)
$$
for every reduced connected closed subscheme $Y \subset X$.
Since $k$ is algebraically closed and $Y$ connected and reduced
we have $H^0(Y, \mathcal{O}_Y) = k$ (Varieties, Lemma
\ref{varieties-lemma-proper-geometrically-reduced-global-sections}).
Hence $\chi(Y, \mathcal{O}_Y) = 1 - \dim H^1(Y, \mathcal{O}_Y)$.
Thus we have to show
$$
3\deg(\omega_X|_Y) > -2 + 2\dim H^1(Y, \mathcal{O}_Y) + \deg(Z \cap Y)
$$
which is true by Lemma \ref{lemma-no-rational-tail}
except possibly if $Y = X$ or if $\deg(\omega_X|_Y) = 0$.
Since $\omega_X$ is ample the second possibility does not
occur (see first lemma cited in this proof). Finally, if
$Y = X$ we can use Riemann-Roch (Lemma \ref{lemma-rr})
and the fact that $g \geq 2$ to see that the inquality holds.
The same argument with $Z = \emptyset$ shows that
$H^1(X, \omega_X^{\otimes 3}) = 0$.
\end{proof}